% Generated draft: requires mathematical and source review.
\documentclass[11pt]{article}
\usepackage[margin=25mm]{geometry}
\usepackage{fontspec}
\setmainfont{Latin Modern Roman}
\usepackage{amsmath,amssymb,mathtools,unicode-math}
\setmathfont{Latin Modern Math}
\usepackage{enumitem,needspace,xurl,hyperref}
\setlength{\parindent}{0pt}
\setlength{\parskip}{5pt}
\setlength{\emergencystretch}{3em}
\newcommand{\N}{\mathbb{N}}
\newcommand{\Z}{\mathbb{Z}}
\newcommand{\Q}{\mathbb{Q}}
\newcommand{\R}{\mathbb{R}}
\newcommand{\C}{\mathbb{C}}
\newcommand{\F}{\mathbb{F}}
\newcommand{\A}{\mathbb{A}}
\newcommand{\PP}{\mathbb{P}}
\newcommand{\E}{\mathbb{E}}
\newcommand{\eps}{\varepsilon}
\newcommand{\dd}{\,\mathrm{d}}
\begin{document}
\section*{UnsolvedMath ID: 800993}
\textbf{Source problem code:} catalogue entry
\subsection{Definitions and mathematical statement}
The supplied question concerns a partial computable function on the reals. The source does not specify the representation of reals, the precise computability convention, or the notion of randomness. Accordingly, the following is a conditional formulation with these choices left explicit, rather than a fully resolved specification.

Let \(X\) be the intended represented space of reals, let \(\lambda\) denote the measure called Lebesgue measure in the source, and let \(\mathcal R\subseteq X\) be the intended class of random reals. For fixed oracle representations, write \(u\leq_T v\) when a code for \(u\) is computable using a code for \(v\) as an oracle, and write
\[
 u<_T v\quad\Longleftrightarrow\quad u\leq_T v\ \text{and}\ v\not\leq_T u.
\]
Under the pointwise interpretation of randomness preservation, a partial function \(f:D\to X\), with \(D\subseteq X\), preserves randomness if
\[
 x\in D\cap\mathcal R\ \Longrightarrow\ f(x)\in\mathcal R.
\]
This interpretation and the choice of \(\mathcal R\) require confirmation from the source context.

\textbf{Question.} Is there a partial computable function \(f:D\to X\) such that
\begin{align*}
 &\forall x,y\in D,\quad f(x)=f(y)\Longrightarrow x=y,\\
 &\forall x\in D\cap\mathcal R,\quad f(x)\in\mathcal R,\\
 &\lambda(D)>0,\\
 &\lambda\bigl(\{x\in D:f(x)<_T x\}\bigr)>0?
\end{align*}
The source describes this as asking for an injective one-way function on the reals; the displayed conditions are its supplied more precise target. No separate computational-complexity definition of one-wayness is supplied.
\subsection{Short English statement}
Does a partial computable injective function on the reals preserve randomness while strictly lowering Turing degree on a set of positive Lebesgue measure, with its domain also having positive measure? The representation of reals, computability convention, and randomness notion remain unspecified in the supplied excerpt. The registry labels the problem open; this draft does not independently verify that status.
\subsection{Sources}
\begin{itemize}
\item Local UnsolvedMath registry, permanent ID \(800993\), exact source code \texttt{catalogue entry}, title \emph{Open Problems in Mathematical Logic --- Problem 5}. The supplied record attributes the problem to Barmpalias, Gao et al., with proposed year 2026, and labels its status \texttt{open}.
\item Barmpalias, Gao et al., \emph{Open Problems in Mathematical Logic} (2026), Problem 5. Supplied source URL: \url{https://arxiv.org/abs/2608.26628}. The background records access on 2026-10-16 and extraction by \texttt{arxiv\_latex}. These are supplied provenance details; the source and its mathematical context have not been independently consulted or verified for this draft.
\end{itemize}
\end{document}
